\vfill\pagebreak
\section{Grouping values: tuples}
\label{sec:collections:tuples}

A function returns one value. Dividing 17 by 5 gives two answers, though: a
quotient of 3 and a remainder of 2. A \textit{tuple} groups a fixed number of
values into a single one, so that the two answers travel together: a function
can return them, a variable can hold them, and another function can receive
them.

\lstinputlisting[style=coloredverbatim, caption={A function that returns two values, \textit{examples/collections/divide.yr}}, label=lst:collections:divide]{examples/collections/divide.yr}
\vspace{-10pt}%
\begin{lstlisting}[style=bashVerb]
(3, 2)
quotient 3, remainder 2
100 = 7 * 14 + 2
\end{lstlisting}

A tuple is written as its values between parentheses, separated by commas, as
\token{(a / b, a \% b)} on line~4. Its type is written the same way, with a type
in place of each value: \token{(i32, i32)}, on line~3, is a pair of
\token{i32}. \token{println} prints a tuple the way it is written, on line~9.

\subsection{Reading the fields}

The values of a tuple are its \textit{fields}, numbered from 0 in the order they
are written. A dot followed by the number of a field reads it: on line~10,
\token{result.0} is the quotient and \token{result.1} the remainder. The number
of fields of a tuple is its \textit{arity}, and is part of its type. A pair has
no field 2, and asking for one is refused when compiling.

\begin{lstlisting}[style=coloredverbatimError, escapechar=@, caption=A field that does not exist]
use std::io;

fn main() {
  let pair = (3, 4);
  println(pair.0 + @\hb{pair.2}@); // error
}
\end{lstlisting}

The compiler says \errmsg{E4214}{tuple access out of bound(2), tuple arity is
  2}.

The fields of a tuple need not have the same type. In the next listing,
\token{person} is a \token{([c8], i32, f64)}: a string, an integer and a
floating point number, which stand for a name, an age and a height.

\lstinputlisting[style=coloredverbatim, caption={A tuple of three types, \textit{examples/collections/person.yr}}, label=lst:collections:person]{examples/collections/person.yr}
\vspace{-10pt}%
\begin{lstlisting}[style=bashVerb]
Alice, 31 years old, 1.68 m tall
Alice: 1.68 m
Alice
31
1.68
\end{lstlisting}

\subsection{Taking a tuple apart}

\token{result.0} and \token{result.1} say where a value is, not what it is. A
tuple can instead be taken apart into variables with names of their own. On
line~12 of Listing~\ref{lst:collections:divide},
\token{let (q, r) = divide(100, 7);} declares two variables at once, one per
field, in order: \token{q} holds the quotient and \token{r} the remainder. The
number of names must be the arity of the tuple.

A field that the program does not need takes the wildcard \token{\_} in place of
a name, as the iterator of a loop does (Section~\ref{sec:control:for}). On
line~7 of Listing~\ref{lst:collections:person}, the age is left out, and only
\token{name} and \token{height} are declared.

\subsection{Passing the fields as arguments: \texttt{expand}}
\label{sec:collections:expand}

A function that takes three parameters cannot receive a tuple of three fields:
the call gives it one argument, where it expects three. The fields could be
passed one by one, as \token{volume(size.0, size.1, size.2)} does on line~9 of
the next listing, but the keyword \token{expand} says it in one word.

\lstinputlisting[style=coloredverbatim, caption={Expanding a tuple, \textit{examples/collections/volume.yr}}, label=lst:collections:volume]{examples/collections/volume.yr}
\vspace{-10pt}%
\begin{lstlisting}[style=bashVerb]
24
24
300
\end{lstlisting}

\token{expand size}, on line~10, replaces the tuple with its fields, separated by
commas, as if they had been written one by one. The call therefore receives
three arguments. The fields need not be all the arguments: on line~13,
\token{expand floor} gives the first two, and \token{10} the third.

Without \token{expand}, \token{volume(floor, 10)} is refused with
\errmsg{E4226}{the call operator is not defined for \errhole{} and
  \errhole{}}, and the note under it says that 3 parameters are expected, where
2 are given.

\token{expand} is not limited to calls. Inside the parentheses of a tuple, it
also stands for the fields of another tuple, which builds a longer tuple from
shorter ones.

\begin{lstlisting}[style=coloredverbatim]
let date = (28, 9, 2026);
let time = (14, 30);
let moment = (expand date, expand time);
println(moment);
let dated = (expand date, "Monday");
println(dated);
\end{lstlisting}
\vspace{-10pt}%
\begin{lstlisting}[style=bashVerb]
(28, 9, 2026, 14, 30)
(28, 9, 2026, Monday)
\end{lstlisting}

On line~3, \token{(expand date, expand time)} is read as if it were
\token{(28, 9, 2026, 14, 30)}: the three fields of \token{date} followed by the
two of \token{time}, a tuple of five fields. On line~5, the fields of
\token{date} are followed by a string, and \token{dated} is an
\token{(i32, i32, i32, [c8])}: the fields keep their types, and the string adds
its own.

\subsection{Going through the fields}

\token{for} goes through the fields of a tuple, from the first to the last, as
on line~10 of Listing~\ref{lst:collections:person}. With two names, the first
holds the number of the field, and the second the field itself.

\begin{lstlisting}[style=coloredverbatim]
let row = ("Alice", 31, 1.68, true);
for i, field in row {
  println(i, ": ", field, ", of type ", typeof(field)::typeid);
}
\end{lstlisting}
\vspace{-10pt}%
\begin{lstlisting}[style=bashVerb]
0: Alice, of type [c8]
1: 31, of type i32
2: 1.68, of type f64
3: true, of type bool
\end{lstlisting}

\token{typeof(field)::typeid} is the name of the type of \token{field}, as a
string; a later chapter comes back to it. It shows what sets this loop apart
from the others: the iterator has the type of the current field, so
\token{field} is a \token{[c8]} in the first iteration, an \token{i32} in the
second, an \token{f64} in the third and a \token{bool} in the last.

The compiler can only do that by copying the block of the loop once for each
field, with a different type for \token{field} in each copy. The number of
iterations is therefore always the arity of the tuple, and the block must be
valid for every type of field. \token{println} is; \token{field + 1} is not,
since a string cannot be added to a number.

\begin{lstlisting}[style=coloredverbatimError, escapechar=@, caption=A block that is not valid for every field]
use std::io;

fn main() {
  let person = ("Alice", 31, 1.68);
  for field in person {
    println(@\hb{field + 1}@); // error
  }
}
\end{lstlisting}

The compiler says \errmsg{E4224}{undefined operator + for types [c8] and i32}:
the copy of the block made for the first field, a string, does not compile,
although the copies for the two numbers would.

\subsection{Comparing tuples}

Two tuples of the same type are equal when their fields are equal, one by one:
\token{divide(17, 5) == (3, 2)} is \token{true}. \token{!=} is its opposite.
Tuples have no order, so \token{<} and \token{>} are not defined on them: there
is no natural way to say which of \token{(3, 2)} and \token{(2, 3)} comes first.

\subsection{A tuple of one value}

Parentheses also group operations, as in \token{(a + b) * 2}, so \token{(42)}
is the number 42, not a tuple. A tuple with a single field is written with a
comma after its value, \token{(42,)}, and its type likewise, \token{(i32,)}.
Such tuples are common in generic code, which the course comes to in later
chapters: a function written for tuples of any arity must work for an arity of
one as well, and a list of arguments collected into a tuple may hold a single
value. They also appear by mistake. When a compiler message names a type such as
\token{(i32,)} where a number was expected, the cause is often a stray comma
after a single value.
